(1) 定向从热端向冷端的总热功率 \(P=-Ak(T)\,dT/dx\) 沿轴恒定。积分 \(\int_0^L(P/A)dx=\int_{T_c}^{T_h}k(T)dT\)，得 \(P=(Ak_0/L)[(T_h-T_c)+\alpha(T_h^2-T_c^2)/2]\)。利用 \(T_h^2-T_c^2=(T_h-T_c)(T_h+T_c)\)，亦可写 \(P=A k_{\rm eff}(T_h-T_c)/L\)，\(k_{\rm eff}=k_0[1+\alpha(T_h+T_c)/2]\)。

(2) 定义严格递增函数 \(F(T)=T+\alpha T^2/2\)（因 \(F'(T)=1+\alpha T>0\)）。在任意位置积分 Fourier 方程，得到 \(F(T(x))=F(T_h)-[F(T_h)-F(T_c)]x/L\)。故 \(T(x)=F^{-1}[F(T_h)-(F(T_h)-F(T_c))x/L]\)，并且 \(dT/dx=-P/[A k(T)]<0\)。当 \(\alpha>0\)，可解显式根 \(T(x)=[-1+\sqrt{1+2\alpha\,F(T(x))}]/\alpha\)，选非负物理支。

(3) 局域熵产生密度 \(\sigma_s=k(T)(T_x)^2/T^2\ge0\)，因为热流从高温流向低温。积分 \(\dot S_{\rm gen}=\int_0^L A k(T)(T_x)^2/T^2dx=P\int_{T_c}^{T_h}dT/T^2=P(1/T_c-1/T_h)>0\)。也可从系统稳态熵账本得到：冷端每秒接受熵 \(P/T_c\)，热端每秒流出 \(P/T_h\)，其差即导热内部不可逆熵产生。

(4) 令 \(\alpha\to0\)，函数 \(F(T)\to T\)，故剖面恢复线性 \(T(x)=T_h-(T_h-T_c)x/L\)，热流 \(P\to A k_0(T_h-T_c)/L\)。弱非线性修正相对于常数热导率的比值精确为 \(1+\alpha(T_h+T_c)/2\)，而非仅使用杆平均温度处的热导率作为不经证明的普适公式；本模型中线性 \(k(T)\) 恰好使该平均表达式精确。
